\section{Álgebra}

\questao{(OM$^3$SP 2026 - 1ª Fase) Em um dia ensolarado Aldemir, coordenador dos Anos Iniciais, propôs um desafio matemático com as letras $M$ e $P$ que compõem a logo da escola. Ele associou essas letras a dois algarismos, do $0$ ao $9$, e escreveu em uma painel da escola que $MP+PM=121$ e $PM-PP=1$. Qual será o valor numérico encontrado para  $M$?}
        \begin{multicols}{4}
                \begin{enumerate}[label=\alph*)]
                    \item 5
                    \item 6
                    \item 7
                    \item 8
                \end{enumerate}
            \end{multicols}

\begin{solucoes}{2}
    \solucao
Nesta questão, tem-se que $M$ e $P$ são algarismos, ou seja, estão entre os números inteiro de 0 a 9. Sabe-se também que  $MP+PM=121$ e $PM-PP=1$. Partindo de $MP+PM=121$, tem-se que, para gerar final 1, as opções são $0$ e $1$, $2$ e $9$, $3$ e $8$, $4$ e $7$ ou $5$ e $6$. Olhando para essas opções, a única que será prontamente descartada é $M=0$ e $P=1$, pois $01+10=11\neq 121$. Todas as outras somadas darão $121$. Com isso, é preciso notar que $PM-PP=1$ e, para isso acontecer, $M-P=1$. O único caso possível seria $M=6$ e $P=5$. Sendo assim, $M=6$.\\
\textbf{ALTERNATIVA B}

\solucao

Podemos pensar de outra maneira e, com isso, podemos resolver de uma forma puramente algébrica. Assim, partindo da primeira informação, 
\begin{equation*}
MP+PM=121
\end{equation*}
\begin{equation*}
    (10M+P)+(10P+M)=121
\end{equation*}
\begin{equation*}
    (10M+M)+(10P+P)=121
\end{equation*}
\begin{equation*}
    11M+11P=121
\end{equation*}
\begin{equation*}
    11*(M+P)=121
\end{equation*}
\begin{equation*}
    M+P=\frac{121}{11}=11
\end{equation*}
\begin{equation*}
    M+P=11 \ (1)
\end{equation*}

Faremos a mesma coisa com a segunda informação,

\begin{equation*}
    PM-PP=1
\end{equation*}
\begin{equation*}
    (10P+M)-(10P+P)=1
\end{equation*}
\begin{equation*}
    M-11M+10M=1
\end{equation*}
\begin{equation*}
    M-P=1 \ (2)
\end{equation*}

Somando $(1)$ e $(2)$,temos que

\begin{equation*}
  2M=12
\end{equation*}
\begin{equation*}
    M=\frac{12}{2}=6
\end{equation*}
Com isso, determinamos que $M=6$. \\
\textbf{ALTERNATIVA B}
\end{solucoes}

\questao{(OM$^3$SP 2026 - 1ª Fase) Na última semana, a Vice-Diretora Quitéria saiu ao portão de entrada da E.M.M.M.S.P. e percebeu que tinham alguns carros e algumas motos estacionadas na frente da escola. Ao entrar ela disse a Félix que o número total de automóveis, apenas carros e motos, presentes na frente da escola era de 12. Ela também contou e determinou que a quantidade de rodas era de 34. Com isso, Félix descobriu quantas motos e quantos carros ali estavam. Qual era o número de motos presentes naquele momento na frenta da E.M.M.M.S.P.?}
        \begin{multicols}{4}
                \begin{enumerate}[label=\alph*)]
                    \item 5
                    \item 6
                    \item 7
                    \item 8
                \end{enumerate}
            \end{multicols}

\begin{solucoes}{2}

\solucao
    Tem-se aqui um clássico problema de sistema de equações do primeiro grau. Tome $C$ representando o número de carros e $M$ o número de motos. Convertendo as informações em equações, ter-se-á:
        \begin{equation*}
\begin{cases}
  C+M=12\\
     4C+2M=34 
\end{cases}
 \end{equation*}
    Ao multiplicar por $-4$ a primeira equação. ter-se-á:
        \begin{equation*}
\begin{cases}
    -4C+-4M=-48\\
     4C+2M=34 
\end{cases}
 \end{equation*}   
Somanda asduasequações acima,teremos
\begin{equation*}
    -2M=-14
\end{equation*} 
\begin{equation*}
    M=\frac{14}{2}=7
\end{equation*}
Assim, na porta da escola havia 7 motos.\\
\textbf{ALTERNATIVA C}

\solucao

Podemos resolver este problema de uma forma menos algébrica e mais aritmética. Pense o seguinte, se todos os veículos fossem motos teríamos $12*2=24$, todavia temos $34$ rodas. Sendo assim, temos 10 rodas a mais do que o número de motos e como cada carro possui o dobro de rodas de cada moto, devemos então dividir $\frac{10}{2}=5$. Logo, serão $5$ carros e assim $7$ motos.  \\
\textbf{ALTERNATIVA C}

\end{solucoes}


